\section{Q: prefix flatness, banked structure, and the exact remaining input}
\label{sec:q-work}

\paragraph{Status.}
This section is an \textup{AUDIT/CONDITIONAL} work snippet for the Q part of
\texttt{experimental/grande\_finale.tex}.  It records the pieces of the
quotient/prefix input that are already theorem-level in the current raw package
\cite{CAP25v13}, and separates them from the still-conjectural max-fiber
flatness needed by the finite adjacent rows.  It does not claim a deployed safe
upper certificate.

\subsection{Definitions and normalizations}

Let \(D\subseteq \B\) be a smooth multiplicative domain of size \(n\).  For
circle/twin-coset rows the same prefix definitions apply after replacing the
multiplicative domain polynomial by the corresponding circle domain polynomial;
the explicit divisor statement below is written only in the multiplicative
case.  Fix an agreement \(m\) and an effective list dimension \(K\), and put
\[
        w:=m-K,\qquad
        \Lambda_M(X):=\prod_{x\in M}(X-x)
        =X^m+\sum_{j=1}^{m}(-1)^j e_j(M)X^{m-j}.
\]
The identity prefix map is
\[
        \Phi_{m,w}:\binom{D}{m}\longrightarrow \B^w,\qquad
        M\longmapsto \bigl((-1)^j e_j(M)\bigr)_{1\le j\le w}.
\]
For \(z\in\B^w\), write
\[
        N_w(z):=\left|\Phi_{m,w}^{-1}(z)\right|,
        \qquad
        \overline F_{m,w}:=\binom{n}{m}|\B|^{-w}.
\]
The number \(\overline F_{m,w}\) is only the ambient average over \(\B^w\); if
one later replaces \(\B^w\) by the actual image of \(\Phi_{m,w}\), the
denominator must be stated explicitly in the row packet.

\begin{proposition}[Newton dictionary]\label{prop:q-newton-dictionary}
If \(\operatorname{char}\B>w\), then the elementary-prefix map
\(\Phi_{m,w}\) is triangularly equivalent to the power-sum map
\[
        \Psi_{m,w}(M):=
        \left(\sum_{x\in M}x^j\right)_{1\le j\le w}.
\]
In particular the two maps have exactly the same fiber sizes.
\end{proposition}

\begin{proof}
Newton's identities give, for \(1\le j\le w\),
\[
        j e_j=\sum_{i=1}^{j}(-1)^{i-1}e_{j-i}p_i,
        \qquad p_i:=\sum_{x\in M}x^i,\quad e_0:=1.
\]
Because every \(j\le w\) is invertible in \(\B\), these identities express
\((e_1,\ldots,e_w)\) and \((p_1,\ldots,p_w)\) as mutually inverse
unitriangular polynomial changes of coordinates.  The sign convention
\((-1)^je_j\) is a coordinatewise bijection, so fiber sizes are unchanged.
\end{proof}

\subsection{What is already proved}

\begin{proposition}[identity prefix floor and exact witness lists]
\label{prop:q-prefix-floor-exact}
For every \(m\ge K\) there is a \(\B\)-valued received word \(U_z\) for
\(\RS[\F,D,K]\) with at least
\[
        \left\lceil \binom{n}{m}|\B|^{-w}\right\rceil
\]
codewords agreeing on \(m\) positions.  More precisely, for every prefix value
\(z\), the list of codewords agreeing with
\[
        U_z(X)=X^m+\sum_{j=1}^{w}z_jX^{m-j}
\]
on at least \(m\) positions is exactly indexed by the fiber
\(\Phi_{m,w}^{-1}(z)\), and no codeword agrees with \(U_z\) on more than \(m\)
positions.
\end{proposition}

\begin{proof}
For \(M\in\Phi_{m,w}^{-1}(z)\), the difference
\[
        U_z-\Lambda_M
\]
has degree at most \(m-w-1=K-1\), hence is a codeword and agrees with \(U_z\)
on \(M\).  Conversely, if a degree-\(<K\) polynomial \(c\) agrees with \(U_z\)
on at least \(m\) points, then \(U_z-c\) is monic of degree \(m\), has those
points as roots, and therefore equals \(\Lambda_M\) for its root set \(M\).
Comparing the top \(w\) coefficients puts \(M\) in the fiber of \(z\), and
then \(c=U_z-\Lambda_M\).  A nonzero degree-\(m\) polynomial cannot vanish on
more than \(m\) points.  This is the content of
\texttt{lem:capff1-identity-prefix-floor} and
\texttt{prop:capff1-witness-exact} in \cite{CAP25v13}.
\end{proof}

\begin{proposition}[head-depth flatness is theorem-level]
\label{prop:q-head-depth}
For the deployed identity-prefix rows, two-sided prefix-fiber flatness is
proved at the following head depths.
\[
\begin{array}{c|c|c}
\text{row family} & \text{deployed subset depths} & \text{head rungs}\\
\hline
\text{KoalaBear multiplicative} & w\le 21 & w\le 22\\
\text{Mersenne-31 circle} & w\le 10 & w\le 11 .
\end{array}
\]
In these ranges every fiber satisfies
\[
        N_w(z)=\binom{n}{m}|\B|^{-w}(1+\varepsilon_z),
        \qquad |\varepsilon_z|\le 2^{-\mathrm{margin}},
\]
with the printed row margins from \texttt{cor:capff1-collision-head} and
\texttt{cor:capff1-circle-head} of \cite{CAP25v13}.  These base cases do not
extend to the frontier depth \(w\approx 6.7\cdot 10^4\).
\end{proposition}

\begin{proof}
This is the Weil/head-depth theorem of the raw package.  For multiplicative
cosets the proof bounds the relevant power-sum character sums after the
power-parametrization of the coset; for circle rows it uses the analogous
additive-character bound on the twin-coset conic.  The error term remains below
the main term exactly in the printed head ranges.  The exact-list statement is
\cref{prop:q-prefix-floor-exact}.
\end{proof}

\begin{proposition}[prefix-collision rigidity]\label{prop:q-prefix-rigidity}
Let \(1\le w<m\le n\).  If \(M\ne M'\) are \(m\)-subsets of \(D\) with
\(\Phi_{m,w}(M)=\Phi_{m,w}(M')\), then
\[
        |M\setminus M'|=|M'\setminus M|\ge w+1.
\]
If equality holds, then
\[
        \Lambda_{M\setminus M'}-\Lambda_{M'\setminus M}
        \in\F^\times ,
\]
so the first nonzero collision stratum is the constant-shift split-pair
stratum.
\end{proposition}

\begin{proof}
Put \(R=M\cap M'\), \(A=\Lambda_{M\setminus M'}\),
\(B=\Lambda_{M'\setminus M}\), \(G=\Lambda_R\), and
\(e=|M\setminus M'|\).  Then
\[
        \Lambda_M-\Lambda_{M'}=G(A-B).
\]
Equality of the monic term and of the depth-\(w\) prefix gives
\(\deg(\Lambda_M-\Lambda_{M'})\le m-w-1\), hence
\(\deg(A-B)\le e-w-1\).  If \(e\le w\), this forces \(A=B\), hence the two
disjoint root sets are equal and \(M=M'\), contradiction.  Thus \(e\ge w+1\).
If \(e=w+1\), then \(A-B\) is a nonzero constant.  This is
\texttt{lem:capg-prefix-rigidity} of \cite{CAP25v13}.
\end{proof}

\begin{corollary}[unconditional anticode cap]\label{cor:q-anticode-cap}
For every \(z\in\B^w\),
\[
        N_w(z)\le {\binom{n}{m-w}\over \binom{m}{w}}.
\]
\end{corollary}

\begin{proof}
By \cref{prop:q-prefix-rigidity}, two distinct members of one fiber intersect
in at most \(m-w-1\) points.  Hence no \((m-w)\)-subset of \(D\) is contained
in two fiber members.  Counting pairs \((R,M)\) with
\(M\in\Phi_{m,w}^{-1}(z)\), \(R\subseteq M\), and \(|R|=m-w\), gives
\[
        N_w(z)\binom{m}{w}\le \binom{n}{m-w}.
\]
\end{proof}

\begin{proposition}[exact second moment]\label{prop:q-second-moment}
Let \(D'\subseteq D\), and for \(e\ge1\) define
\[
\operatorname{sp}_w(e;D') :=
\#\left\{(A,B):
\begin{array}{l}
\text{\(A,B\) monic split over \(D'\), degree \(e\),}\\
\text{their root sets are disjoint, and } \deg(A-B)\le e-w-1
\end{array}\right\}.
\]
Then
\[
        \sum_{z\in\B^w}N_w(z)^2
        =
        \binom{n}{m}
        +\sum_{e=w+1}^{\min(m,n-m)}
          \sum_{\substack{R\subseteq D\\ |R|=m-e}}
          \operatorname{sp}_w(e;D\setminus R).
\]
Moreover the top stratum is
\[
\operatorname{sp}_w(w+1;D')=
\#\{(A,c):\ A\text{ monic split over }D'\text{ of degree }w+1,
c\in\F^\times,\ A-c\text{ split over }D'\}.
\]
\end{proposition}

\begin{proof}
The left side counts ordered pairs \((M,M')\) in a common prefix fiber.  The
diagonal contributes \(\binom{n}{m}\).  For \(M\ne M'\), write
\[
        R=M\cap M',\qquad
        A=\Lambda_{M\setminus M'},\qquad
        B=\Lambda_{M'\setminus M}.
\]
This is a bijection between non-diagonal colliding pairs and triples with
\(|R|=m-e\), \(A,B\) split of degree \(e\) over \(D\setminus R\), disjoint root
sets, and \(\deg(A-B)\le e-w-1\).  The range \(e\ge w+1\) is
\cref{prop:q-prefix-rigidity}; also \(e\le m,n-m\).  At \(e=w+1\), the degree
condition says \(A-B=c\in\F^\times\), and conversely such a constant shift has
disjoint roots because a common root would force \(c=0\).  This is the exact
ledger \texttt{thm:capg-second-moment} in \cite{CAP25v13}.
\end{proof}

\begin{proposition}[twist-orbit constancy]\label{prop:q-orbit}
For a multiplicative domain stable under \(\mu_n\), the action
\[
        (z_1,\ldots,z_w)\longmapsto(\zeta z_1,\zeta^2z_2,\ldots,\zeta^wz_w),
        \qquad \zeta\in\mu_n,
\]
preserves fiber sizes.
\end{proposition}

\begin{proof}
The map \(M\mapsto \zeta^{-1}M\) is a bijection on \(m\)-subsets of \(D\), and
\[
        \Lambda_{\zeta^{-1}M}(X)=\zeta^{-m}\Lambda_M(\zeta X).
\]
Thus the \(j\)-th prefix coordinate is multiplied by \(\zeta^j\).  This is
\texttt{lem:capfp-orbit} of \cite{CAP25v13}.
\end{proof}

\begin{proposition}[boundary divisor form for multiplicative rows]
\label{prop:q-boundary-divisor}
Assume \(D\) is a multiplicative coset with
\(\Lambda_D=X^n-\beta\).  For \(z\in\B^w\), set
\[
        P_z(X):=X^m+\sum_{j=1}^{w}z_jX^{m-j},
        \qquad X^n=P_zQ_z+R_z,\quad \deg R_z<m,
\]
and \(\omega:=n-m\).  In the boundary branch of the interpolation lattice,
the fiber of \(z\) is exactly the prescribed-top-coefficient divisor census
\[
        N_w(z)
        =
        \#\left\{A:\ \deg A\le \omega-w-1,\quad Q_z+A\mid X^n-\beta\right\}.
\]
Equivalently, Q is the boundary profile \(d_1=w+1\) of the split-pencil
problem, while the BC input begins at the interior profiles.
\end{proposition}

\begin{proof}
The pair \((1,P_z)\) belongs to the interpolation lattice of \(U_z\) and has
shifted degree \(w+1\).  The Euclidean division gives
\[
        (Q_z,\beta-R_z)
\]
as a second lattice vector, since at every \(x\in D\),
\[
        Q_z(x)P_z(x)=x^n-R_z(x)=\beta-R_z(x).
\]
Their determinant is \(\beta-X^n=-\Lambda_D\), so they form a lattice basis.
In the branch where the minimal shifted degree is exactly \(w+1\), every census
element of \(W\)-degree \(\omega\) is uniquely \(Q_z+A\) after monic
normalization, with \(\deg A\le\omega-w-1\).  The split condition is exactly
\((Q_z+A)\mid X^n-\beta\).  This is the statement
\texttt{thm:capfp-Q-unify} in \cite{CAP25v13}.
\end{proof}

\subsection{The exact remaining Q input}

\begin{conjecture}[Q max-fiber flatness, active form]\label{conj:q-active}
After quotient-pulled-back, planted, and other already paid prefix strata have
been charged at their own image scale, the primitive identity-prefix fibers
satisfy
\[
        N_w^{\mathrm{prim}}(z)
        \le
        \kappa_Q(n)\binom{n}{m}|\B|^{-w}
        \qquad\text{for all }z\in\B^w.
\]
For asymptotic frontier closure with agreement reserve \(R\gg\log n\), it is
enough to prove \(\kappa_Q(n)=n^{O(1)}\), or more generally
\(\kappa_Q(n)=e^{o(n)}\).  For the finite adjacent rows, \(\kappa_Q\) must be
an explicit row-dependent constant and the first-match quotient ledger must fit
inside the printed margins \(22.2,22.0,3.3,3.1\) bits after all other paid
cells are included.
\end{conjecture}

\begin{remark}[what is not proved by the current Q lemmas]
\label{rem:q-not-proved}
\Cref{cor:q-anticode-cap} is a genuine all-depth theorem, but it remains far
above the average fiber at the frontier depths.  \Cref{prop:q-second-moment}
is an exact identity, but it is not a max-fiber bound until the primitive
shift-pair terms are bounded and higher moments or an extremality theorem
control worst-case fibers.  In particular, second-moment control alone cannot
settle the printed adjacent rows.
\end{remark}

\begin{proposition}[moment sandwich]\label{prop:q-moment-sandwich}
Let \(\mu(z):=N_w(z)/\binom{n}{m}\), define
\[
        \Gamma_r:=|\B|^{w(r-1)}\sum_z\mu(z)^r,\qquad r\ge2,
\]
and let \(R:=|\B|^w\max_z\mu(z)\) be the max-to-mean ratio.  Then
\[
        { \log_2\Gamma_r\over r-1}
        \le
        \log_2 R
        \le
        { \log_2\Gamma_r+w\log_2|\B|\over r}.
\]
Thus \(R\le e^{o(n)}\) follows from \(\Gamma_r\le e^{o(n)}\) for every fixed
\(r\), but a direct moment-to-max proof at the deployed adjacent margins would
need a row-dependent finite hierarchy of very high order.
\end{proposition}

\begin{proof}
Since \(N_w(z)\le R\binom{n}{m}|\B|^{-w}\), one has
\(\Gamma_r\le R^{r-1}\), giving the left inequality.  Conversely,
\[
        \max_z\mu(z)\le \left(\sum_z\mu(z)^r\right)^{1/r}
        =\left(\Gamma_r|\B|^{-w(r-1)}\right)^{1/r},
\]
and multiplying by \(|\B|^w\) gives the right inequality.  If every fixed
\(r\) has \(\Gamma_r=e^{o(n)}\) and \(w\log_2|\B|=O(n)\), choose \(r\) as a
large constant after fixing the desired \(o(n)\) slope.  This reproduces
\texttt{prop:capfp-Q-sandwich} of \cite{CAP25v13}.
\end{proof}

\subsection{Proof strategy that is still legitimate}

\begin{enumerate}[label=\textup{Q\arabic*.},leftmargin=3.0em]
\item \textbf{Run the finite rung audit before analytic work.}
At \(a_0+1\), quotient-periodic support counts are too large at support level
and must be paid by image descent and first-match coalescing.  The audit must
print every divisor rung, image denominator, loss in bits, and residual
primitive bucket.  A tight or inverted rung is a periodic obstruction, not a
primitive Q failure.

\item \textbf{Use the proved collision ledger to build fixed moments.}
\Cref{prop:q-second-moment} gives the exact \(r=2\) object.  Higher fixed
moments should be organized as collision hypergraphs of split locators with
common top prefix.  The asymptotic Q target is equivalent, via
\cref{prop:q-moment-sandwich}, to proving \(\Gamma_r=e^{o(n)}\) for every
fixed \(r\), after quotient-pulled-back cells have been removed.

\item \textbf{Exploit twist symmetry to force structure on heavy fibers.}
\Cref{prop:q-orbit} makes fiber size constant on twist orbits.  Moment bounds
make heavy fibers few; if there are fewer heavy values than a generic orbit,
then a heavy \(z\) must have nontrivial stabilizer, forcing its nonzero
coordinates into divisible index classes.  That is exactly the quotient-pulled
back shape that should descend to a coarser rung, modulo explicitly tracked
folding defects.

\item \textbf{Reduce the finite problem to extremality.}
The finite adjacent constants are too small for unspecified polynomial losses.
The clean finite target remains mode-at-null or exchange compression:
\[
        N_w(z)\le N_w(0)\quad\text{after quotient rungs are charged},
        \qquad
        N_w(0)\le \kappa_{\mathrm{row}}\overline F_{m,w}.
\]
This is still conjectural.  It is finitely falsifiable, and a violation would
identify a new primitive prefix obstruction floor.

\item \textbf{Do not use character-sum head-depth results beyond their range.}
The Weil/head-depth proof pays one square-root cost per active prefix coordinate
and stops at \(w\le21\)--\(22\) for KoalaBear and \(w\le10\)--\(11\) for
Mersenne-31.  The frontier depth is about \(6.7\cdot10^4\).  Per-frequency
short-interval or monodromy estimates, even if individually sharp, do not by
themselves prove Q; the missing cancellation is joint cancellation across the
whole character-tuple family, as recorded in
\texttt{prop:capfp-Q-smooth} and \texttt{rem:capfp-Q-wall} of
\cite{CAP25v13}.
\end{enumerate}

\paragraph{Overclaim guard.}
The statements above prove structural control of Q and exact calibration of its
collision ledger.  They do not prove \(\cref{conj:q-active}\), do not prove the
first safe adjacent rows, and do not authorize spending a challenge-field
denominator on a base-field prefix fiber.  Any promotion into
\texttt{experimental/grande\_finale.tex} should keep the final Q input visibly
CONJECTURAL unless a row packet supplies the finite constants or a theorem
supplies the asymptotic max-fiber flatness.
